\chapter{Evaluation}
Following the project proposal, the evaluation of \textit{Coherence} focuses on correctness and expressive power.

\section{Correctness}
There is a wide range of potential inaccuracies in this project. The compilation pipeline might contain bugs, the runtime could suffer from heisenbugs, and the type system may be unsound. This section discusses the methods used to rule out these issues.

\subsection{Testing}
% (92 words)
To detect bugs in the compilation pipeline implementation, a test suite of 62 tests was developed. The largest portion of this suite consists of 47 type-checking tests. They were designed to cover a wide range of programming constructs and target distinct components of the type-checking stage. The suite includes both positive and negative tests to ensure that the compiler accepts valid code and rejects incorrect code. Additionally, 4 unit tests were created to verify the correct implementation of lock-inference algorithm. The remaining 11 tests perform end-to-end verification of the entire compilation pipeline (from parsing to codegen and linking with the runtime). Of these, six are small, deterministic programs intended to confirm that core language constructs in Coherence compile as expected and interact properly with the runtime. The other five are concurrency-focused. These tests target potential runtime bugs and were also used to validate claims of deadlock freedom. Most of these tests were large in scale, involving the creation of numerous actor instances and the transmission of many messages. This approach was necessary because heisenbugs were difficult to detect in short programs, and many bugs only became apparent when the tests were scaled up.


\subsection{Programming principles}
% 204 words
In addition to testing, to prevent bugs, the compiler extensively utilizes defensive programming. Assertions are integrated throughout the compilation pipeline to verify internal invariants. These checks act as an immediate signal for logic bugs.\\
Implementing the runtime correctly proved challenging due to frequent encounters with edge-case heisenbugs. To ensure correctness beyond standard testing, a structured programming paradigm was introduced. Shared runtime data structures were protected with locks to prevent data races. However, this alone was insufficient to guarantee correct behaviour.  Many operations needed to be performed atomically---for example, adding an actor instance to a lock's wait queue while simultaneously changing its state to WAITING. To address this, all such atomic operations were identified and wrapped in logical atomic sections. These sections were implemented by acquiring all necessary locks at the start of the section and releasing them at the end. The runtime was designed to always acquire locks in a fixed global order to prevent deadlocks. These principles closely mirror the strategies used by the \textit{Coherence} language to prevent concurrency bugs. While this creates a natural application of these concepts, implementing the runtime directly in \textit{Coherence} is currently infeasible due to the lack of access to low-level OS constructs like threads and user-level context switches.


\subsection{Correctness of type-system}
% 647 words
In the context of \textit{Coherence}, soundness is defined to be the absence of data races and deadlocks.

Deadlock in \textit{Coherence} is defined as a state where the runtime's schedule queue is empty and no actor instances are runnable, yet some actor instances have work left (their mailboxes are not empty). This can only occur when actor instances enter a cyclic wait within lock queues. This is prevented by \textit{Coherence} as the generated code takes locks in a globally fixed order, preventing cyclic dependencies. 

For data race safety, we first address \textit{Coherence1.0} which is much simpler as it does not have to deal with the complexities of viewpoint adaptation. Since \texttt{ref}, \texttt{val}, and \texttt{iso} are inherited directly from Pony, their data race freedom is already formally established. The soundness argument for \textit{Coherence1.0} therefore focuses on the \texttt{locked<L>} capability and how it interacts with existing capabilities. The argument relies on two lemmas. The first states that if all reference capability invariants hold, no sequence of memory accesses can result in a data race. For \texttt{locked<L>}, the invariant requires that any read or write to the reference can only occur when the accessing actor holds lock \texttt{L}. Since mutual exclusion ensures only one actor can hold \texttt{L} at a time, all accesses are synchronized. The second lemma establishes that all execution steps (aliasing, message passing, and so on) preserve these invariants. This holds because \texttt{locked<L>} can only alias with references of the same capability, and assignment from a different capability is only permitted when the source has no existing aliases. Together, these lemmas show that invariants hold at every execution step, and that invariant preservation implies safety throughout execution.

Establishing data-race safety for \textit{Coherence 2.0} is considerably more complex. The interactions involving \texttt{locked} capabilities are more nuanced—for instance, we must now account for \texttt{locked<L>} appearing as a field type within nested pointer structures. Additionally, as discussed in section \ref{...}, the type checking in \textit{Coherence 2.0} differs subtly from Pony. Fortunately, the assignability rules and the definitions of covariance and invariance are designed to ensure that assignments within a local scope do not violate any reference capability invariant, regardless of the specific definition of $\triangleright$. 

However, these rules are not enough to guarantee data race safety for an arbritary $\triangleright$. Sendability is verified by checking that the outer capability is not \texttt{ref}, but this relies on $\triangleright$ being appropriately defined. Consider, for instance, if $\text{val} \triangleright \text{ref} = \text{ref}$. A behaviour could then accept an \texttt{(int ref) val} as an argument, allowing actors to wrap a \texttt{ref} inside a \texttt{val} and share it, thereby violating the invariant of \texttt{ref}. Similarly, these rules do not guarantee safety when unaliasing an \texttt{iso}. For example, suppose $\text{iso} \triangleright \text{ref} = \text{tag}$ while $\text{iso} \ \triangleright \text{ref} = \text{ref}$. This would permit storing a \texttt{tag} to a \texttt{ref} field of an \texttt{iso}, then unaliasing to obtain a \texttt{ref} to what was originally a \texttt{tag}. This effectively allows assignment of a \texttt{tag} to a \texttt{ref}. Since any reference capability can be assigned to \texttt{tag}, this would permit aliasing arbitrary capabilities with \texttt{ref}, potentially violating the invariants of those capabilities. Hence, for safety, we require that $\text{iso\textasciicircum} \ \triangleright\  \kappa$ is covariant with $\text{iso} \triangleright \kappa$ for all capabilities $\kappa$ when $\text{iso} \ \triangleright \ \kappa$ is not iso. If $\text{iso} \ \triangleright \ \kappa$ is iso, we have the flexibility to set $\text{iso\textasciicircum} \ \triangleright \ \kappa$ to be iso\textasciicircum. 


Finally, we must verify that the definition of $\triangleright$ satisfies the property mentioned in section \ref{...}: when base types are identical, assignability, covariance, and invariance reduce to checking only the outer capabilities.

If the $\triangleright$ operator is associative, the above claims (which are quantified over all types with unbounded pointer nesting) reduce to checking only two levels of indirection (a pointer to a pointer to a base type). This makes exhaustive verification feasible. To this end, I wrote a prolog program that verifies associativity of $\triangleright$, soundness of sendability, soundness of \texttt{iso} unaliasing, and the property claimed in Section~\ref{...}. This program can be found in section (...).

\section{Expressive power}
Perhaps can talk about accessing shared data structures (implementation in Pony and kappa). And the corresponding code in pony (take from the preparation).
Read about rust.
A talk about reference capabilities